\documentclass{beamer}

% Theme selection
\usetheme{Madrid}
\usecolortheme{default}

% Packages
\usepackage{listings} % For syntax highlighting of Python code
\usepackage{xcolor}
\usepackage{hyperref}
\usepackage{graphicx}

% Define Python code style for listings
\definecolor{codegreen}{rgb}{0,0.6,0}
\definecolor{codegray}{rgb}{0.5,0.5,0.5}
\definecolor{codepurple}{rgb}{0.58,0,0.82}
\definecolor{backcolour}{rgb}{0.95,0.95,0.92}

\lstdefinestyle{mystyle}{
    backgroundcolor=\color{backcolour},   
    commentstyle=\color{codegreen},
    keywordstyle=\color{magenta},
    numberstyle=\tiny\color{codegray},
    stringstyle=\color{codepurple},
    basicstyle=\ttfamily\footnotesize,
    breakatwhitespace=false,         
    breaklines=true,                 
    captionpos=b,                    
    keepspaces=true,                 
    numbers=none,                    
    showspaces=false,                
    showstringspaces=false,
    showtabs=false,                  
    tabsize=2
}
\lstset{style=mystyle, language=Python}

% Title page details
\title{Introduction to JAX}
\subtitle{A High-Performance Numerical Computing Library}
\author{Based on learn-jax}
\date{\today}

\begin{document}

% Title Slide
\begin{frame}
    \titlepage
\end{frame}

% Outline Slide
\begin{frame}{Outline}
    \tableofcontents
\end{frame}

\section{What is JAX?}
\begin{frame}{What is JAX?}
    JAX is Autograd and XLA, brought together for high-performance machine learning research.
    
    \vspace{0.5cm}
    \textbf{Core Features:}
    \begin{itemize}
        \item \textbf{NumPy Drop-in:} Provides a familiar \texttt{jax.numpy} API.
        \item \textbf{Hardware Acceleration:} Native, seamless support for CPU, GPU, and TPU.
        \item \textbf{Composable Transformations:} Powerful functions like \texttt{jit}, \texttt{vmap}, and \texttt{grad}.
    \end{itemize}
\end{frame}

\section{Just-In-Time Compilation}
\begin{frame}[fragile]{Transformation 1: \texttt{jax.jit}}
    \textbf{Just-In-Time Compilation} compiles your Python code into optimized machine code using XLA (Accelerated Linear Algebra).

    \vspace{0.3cm}
    \begin{lstlisting}
import jax.numpy as jnp
from jax import jit

def slow_f(x):
    # Standard Python function
    return jnp.dot(x, x.T)

# Fast, compiled version
fast_f = jit(slow_f)
    \end{lstlisting}
    
    \vspace{0.3cm}
    \textit{Benefit: Fuses operations, dramatically reducing memory overhead and execution time.}
\end{frame}

\section{Automatic Vectorization}
\begin{frame}[fragile]{Transformation 2: \texttt{jax.vmap}}
    \textbf{Automatic Vectorization} transforms a function that operates on a single data point into one that operates efficiently on batches of data.

    \vspace{0.3cm}
    \begin{lstlisting}
from jax import vmap
import jax.numpy as jnp

# Function for a single example
def predict(weights, inputs):
    return jnp.dot(weights, inputs)

# Vectorized to handle batches of inputs
batch_predict = vmap(predict, in_axes=(None, 0))
    \end{lstlisting}
    
    \vspace{0.3cm}
    \textit{Benefit: Eliminates complex, slow manual batching loops while maintaining readability.}
\end{frame}

\section{Automatic Differentiation}
\begin{frame}[fragile]{Transformation 3: \texttt{jax.grad}}
    \textbf{Automatic Differentiation} computes the exact gradient (derivative) of a scalar-valued function automatically.

    \vspace{0.3cm}
    \begin{lstlisting}
from jax import grad

def f(x):
    return 3 * x**2 + 2 * x + 1

# df/dx = 6x + 2
df_dx = grad(f)

print(df_dx(2.0)) # Output: 14.0
    \end{lstlisting}
    
    \vspace{0.3cm}
    \textit{Benefit: Provides exact analytical derivatives without the need for manual calculus.}
\end{frame}

\section{Pytrees}
\begin{frame}[fragile]{Working with Pytrees}
    JAX treats nested Python structures (lists, tuples, dicts) as \textbf{Pytrees}. You can use \texttt{jax.tree\_map} to apply functions across these structures effortlessly.

    \vspace{0.3cm}
    \begin{lstlisting}
from jax.tree_util import tree_map

# A nested dictionary (Pytree)
params = {'w': [1.0, 2.0], 'b': 0.5}

# Double everything in the dictionary
doubled = tree_map(lambda x: x * 2, params)
    \end{lstlisting}
    
    \vspace{0.3cm}
    \textit{Benefit: Easily update complex model parameters without writing custom traversal code.}
\end{frame}

\section{Summary}
\begin{frame}{Summary}
    JAX gives you the power to write standard Python/NumPy code and instantly make it:
    
    \vspace{0.3cm}
    \begin{enumerate}
        \item \textbf{Fast} via \texttt{jax.jit}
        \item \textbf{Batched} via \texttt{jax.vmap}
        \item \textbf{Differentiable} via \texttt{jax.grad}
    \end{enumerate}

    \vspace{0.5cm}
    \begin{center}
        \Large \textbf{Thank You!}
    \end{center}
\end{frame}

\end{document}
